\pdfvariable trailerid {[<C3312026090300000000000000000000><C3312026090300000000000000000000>]}
\documentclass[10pt]{article}
\usepackage[margin=0.78in]{geometry}
\usepackage{amsmath,amssymb,amsthm,microtype,hyperref}
\hypersetup{colorlinks=true,linkcolor=blue,urlcolor=blue}
\providecommand{\CRevisionRound}{2}
\newtheorem{theorem}{Theorem}
\newtheorem{remark}[theorem]{Remark}
\newcommand{\R}{\mathbb R}
\setlength{\emergencystretch}{4em}
\title{Dirac-Monopole Dynamics on the Sphere:\\Magnetic Circles and the Complete Covariant-Laplacian Spectrum}
\author{Route-A source-local certificate HCS-C331}
\date{3 September 2026}

\begin{document}
\maketitle
\begin{abstract}
We identify a conserved Poincar\'e vector and every magnetic circle on the
unit sphere, including its least period and all zero-energy and zero-charge
faces.  This is the \emph{Poincare vector and every magnetic circle} stage.
\ifnum\CRevisionRound>0
For an integral charge we then derive the \emph{complete monopole-harmonic
spectrum}, its multiplicities, heat trace, and charge-conjugation symmetry.
\fi
\ifnum\CRevisionRound>1
Finally we record the \emph{Route-A obstruction and finite certificate}:
Chern integrality and natural quantization do not create a rational-prime
owner, an orbit determinant, or a Hilbert--P\'olya operator.
\fi
\end{abstract}

\section{The classical flow}
Let $S^2\subset\R^3$ be the oriented unit sphere and let
$Jv=x\times v$ on $T_xS^2$.  Fix $q\in\mathbb Z$, put $b=q/2$, and consider
\begin{equation}\label{eq:magnetic}
 \nabla_t\dot x=bJ\dot x,\qquad |\dot x|^2=2E.
\end{equation}

\begin{theorem}[All magnetic trajectories]\label{thm:classical}
The vector
\[
 K=x\times\dot x+b x
\]
is constant and satisfies
\[
 K\cdot x=b,\qquad |K|^2=2E+b^2.
\]
If $E>0$, the trajectory is the oriented small circle cut out by
$K\cdot x=b$ and its least positive period is
\begin{equation}\label{eq:period}
 T(E,q)=\frac{2\pi}{\sqrt{2E+q^2/4}}.
\end{equation}
If $E=0$, the trajectory is stationary.  If $q=0$ and $E>0$, the circle is
a great circle.
\end{theorem}

\begin{proof}
Skew-symmetry of $J$ preserves $|\dot x|$.  The ambient acceleration is
$\ddot x=b\,x\times\dot x-2E x$.  Consequently,
\[
 \dot K=x\times\ddot x+b\dot x
 =b\,x\times(x\times\dot x)+b\dot x=0.
\]
Orthogonality gives the two displayed identities.  The triple-product
identity also gives $K\times x=\dot x$, so $x$ is a rigid rotation about the
fixed axis $K$ with angular speed $|K|$.  Its squared circle radius is
$2E/(2E+b^2)$, positive exactly when $E>0$.  A nonconstant circle first
returns with its tangent after angle $2\pi$, proving~\eqref{eq:period}.
The two boundary statements follow directly.
\end{proof}

Time reversal gives the precise sign pairing: if $x_q(t)$ solves charge $q$,
then $x_q(-t)$, with its initial velocity reversed, solves charge $-q$ and
traverses the same geometric circle oppositely.  This does not identify the
solutions having the same initial velocity.  The conserved-plane proof does
not rely on coordinates or a gauge potential.

\ifnum\CRevisionRound>0
\section{Integral charge and the complete quantum spectrum}
Let $L_q\to S^2$ be a Hermitian line bundle with unitary connection of
curvature
\begin{equation}\label{eq:curvature}
 F_\nabla=i\frac q2\,dA.
\end{equation}
Its first Chern number is
\[
 \frac{1}{2\pi i}\int_{S^2}F_\nabla=q.
\]
Thus integrality of $q$ is exactly the global bundle boundary; a
nonintegral local magnetic charge is not admitted to this quantum model.

Complex line bundles over $S^2$ are classified by
$c_1\in H^2(S^2;\mathbb Z)\cong\mathbb Z$.  Hence $L_q$ is unitarily
isomorphic to the standard degree-$q$ homogeneous monopole bundle.  After
transport to one bundle, two unitary connections with curvature
\eqref{eq:curvature} differ by $i\alpha$, where $d\alpha=0$.  Since
$H^1_{\mathrm{dR}}(S^2)=0$, $\alpha=df$; the gauge $g=e^{if}$ carries one
connection to the other.  Thus the arbitrary frozen connection is
unitarily gauge equivalent to the standard homogeneous one.

On $C^\infty(S^2,L_q)=C_c^\infty(S^2,L_q)$ set
\[
 \mathfrak q_q[s]=\int_{S^2}|\nabla s|^2\,dA,
\]
and let $\Delta_q$ be its Friedrichs realization.  The symmetric
Laplace-type expression $\nabla^*\nabla$ on a compact manifold without
boundary is essentially self-adjoint on smooth sections.  Its unique
self-adjoint closure therefore equals this positive Friedrichs operator;
the resolvent is compact.  Gauge-equivalent connections give unitarily
conjugate operators.
\begin{theorem}[All monopole-harmonic levels]\label{thm:spectrum}
The complete spectrum is
\begin{equation}\label{eq:spectrum}
 \lambda_{n,q}=n(n+|q|+1)+\frac{|q|}{2},\qquad n=0,1,\ldots,
\end{equation}
and the multiplicity of $\lambda_{n,q}$ is $2n+|q|+1$.  Hence, for $t>0$,
\begin{equation}\label{eq:heat}
 \operatorname{Tr}e^{-t\Delta_q}
 =\sum_{n=0}^{\infty}(2n+|q|+1)e^{-t\lambda_{n,q}}.
\end{equation}
The degree $-q$ bundle is the complex conjugate of $L_q$ and has the same
spectrum.
\end{theorem}

\begin{proof}
It now suffices to use the standard homogeneous connection.  Its sections
are degree-$q$ equivariant functions on $SU(2)$.  Peter--Weyl
decomposition contains exactly one copy of each spin
$\ell=|q|/2+n$, $n\geq0$.  The horizontal Laplacian is the total Casimir
minus the fixed vertical weight square; it therefore acts by
\[
 \ell(\ell+1)-\frac{q^2}{4}
 =n(n+|q|+1)+\frac{|q|}{2}.
\]
The spin space has dimension $2\ell+1=2n+|q|+1$.  Peter--Weyl completeness
excludes further levels.  The operator realization above and quadratic
eigenvalue growth justify the absolutely convergent heat trace.  Complex
conjugation changes $q$ to $-q$, while~\eqref{eq:spectrum} depends only on
$|q|$.
\end{proof}

At $q=0$, formula~\eqref{eq:spectrum} becomes the ordinary spherical
harmonic spectrum $n(n+1)$ with multiplicity $2n+1$.  At $n=0$, the lowest
monopole level has energy $|q|/2$ and multiplicity $|q|+1$.
\fi

\ifnum\CRevisionRound>1
\section{Certificate and obstruction}
The checked artifact contains 39 positive-energy classical rows, 195
quarter-turn samples, 357 exact spectral cells, 36 degree-80 heat-trace
partial sums, 25 Chern rows, and 18 time-reversal pairings, for 3,167
audited scalar leaves.  An independent implementation performs 4,414
checks.  SymPy verifies 2,621 exact identities, two isolated reproductions
are byte-identical, and 80
hostile parser, schema, semantic, and repaired-hash attacks are rejected.
These finite data are convention-sensitive regression receipts; the proofs
of Theorems~\ref{thm:classical} and~\ref{thm:spectrum} establish the
all-parameter statements.

The nearest registered owners are C313, the $q=0$ round-sphere boundary;
C289, noncompact hyperbolic magnetic flow; C293, singular magnetic Grushin
dynamics; and C274, the Euclidean Penning trap.  None owns the simultaneous
Chern, magnetic-circle, and monopole-spectrum contract.

The conservative Route-A tuple is
\begin{gather*}
 (\mathrm{A0\_WEAK\_ARITHMETIC\_RELATION},\mathrm{A1\_WEAK},\\
 \mathrm{A2\_FAIL},\mathrm{A3\_FAIL},
 \mathrm{A4\_NATURAL\_QUANTIZATION}).
\end{gather*}
Chern integrality is not a rational-prime or prime-power carrier.  The
circles form clean continuous families rather than isolated hyperbolic
cycles.  The heat trace is not an Euler product or target determinant, and
the covariant Laplacian is not asserted to realize target zeros.  Route A is
therefore rejected and Route B remains locked under
\texttt{NO\_BAD\_EULER\_OR\_ROOT\_NUMBER}.  We claim no target local data,
root number, automorphy, divisor, functional equation, zero match, or
Hilbert--P\'olya operator.

\paragraph{AI use.}
A generative language model assisted drafting and code scaffolding.  The
displayed proofs, independent recomputation, hostile tests, and deterministic
artifacts define the audit record.
\fi

\section*{Source lineage}
This paper is a source-local reconstruction and makes no priority claim.
\begin{thebibliography}{9}
\bibitem{Dirac}
P. A. M. Dirac, ``Quantised Singularities in the Electromagnetic Field,''
\emph{Proc. Roy. Soc. A} 133 (1931), 60--72. DOI:
\href{https://doi.org/10.1098/rspa.1931.0130}{10.1098/rspa.1931.0130}.
\bibitem{WuYang1976}
T. T. Wu and C. N. Yang, ``Dirac Monopole Without Strings: Monopole
Harmonics,'' \emph{Nucl. Phys. B} 107 (1976), 365--380. DOI:
\href{https://doi.org/10.1016/0550-3213(76)90143-7}%
{10.1016/0550-3213(76)90143-7}.
\bibitem{WuYang1977}
T. T. Wu and C. N. Yang, ``Some Properties of Monopole Harmonics,''
\emph{Phys. Rev. D} 16 (1977), 1018--1021. DOI:
\href{https://doi.org/10.1103/PhysRevD.16.1018}%
{10.1103/PhysRevD.16.1018}.
\bibitem{WuYangClassical}
T. T. Wu and C. N. Yang, ``Dirac's Monopole Without Strings: Classical
Lagrangian Theory,'' \emph{Phys. Rev. D} 14 (1976), 437--445. DOI:
\href{https://doi.org/10.1103/PhysRevD.14.437}%
{10.1103/PhysRevD.14.437}.
\end{thebibliography}
\end{document}
